% HFH-LaTeX-Vorlage fuer wissenschaftliche Arbeiten
% Copyright (c) 2026 Ilya Zarubin
%
% Diese Vorlage ersetzt nicht die jeweils aktuellen offiziellen Regelungen.
% Sie erhebt keinen Anspruch auf Vollstaendigkeit oder Fehlerfreiheit. Vor
% jeder Verwendung sind die aktuellen und studiengangsspezifischen
% HFH-Vorgaben fuer Seminar-, Haus-, Bachelor- und Masterarbeiten zu pruefen.
% Bei Abweichungen gelten die offiziellen Vorgaben der HFH, des Pruefungsamts
% und der betreuenden beziehungsweise pruefenden Personen.
%
% Diese Datei darf unter den Bedingungen der LaTeX Project Public License,
% Version 1.3c oder jeder spaeteren Version, verbreitet und veraendert werden.
% Einzelheiten stehen in der Datei LICENSE.
%
% Kompilieren mit pdfLaTeX und BibTeX, zum Beispiel: latexmk -pdf main.tex
\documentclass[
  12pt,
  a4paper,
  oneside,
  headings=normal,
  numbers=noendperiod,
  listof=totoc,
  bibliography=totoc
]{scrreprt}

% ---------------------------------------------------------------------------
% Metadaten: Diese Werte vor dem Schreiben anpassen.
% ---------------------------------------------------------------------------
\newcommand{\MainLanguage}{ngerman} % Fuer Englisch "english" einsetzen.
\newcommand{\ThesisType}{Bachelorarbeit}
\newcommand{\ThesisTitle}{Titel der wissenschaftlichen Arbeit}
\newcommand{\ThesisSubtitle}{Optionaler Untertitel}
\newcommand{\AuthorName}{Vorname Nachname}
\newcommand{\MatriculationNumber}{12345678}
\newcommand{\DegreeProgramme}{Bezeichnung des Studiengangs}
\newcommand{\Department}{Fachbereich}
\newcommand{\FirstExaminer}{Titel Vorname Nachname}
\newcommand{\SecondExaminer}{Titel Vorname Nachname}
\newcommand{\SubmissionDate}{TT. Monat JJJJ}

% Der Abstract bleibt sichtbar, damit seine Einschraenkung nicht uebersehen
% wird. Fuer Bachelor- und Seminararbeiten den Schalter auf false setzen.
\newif\ifIncludeAbstract
\IncludeAbstracttrue

% Optionale Verzeichnisse nur bei Bedarf einschalten.
\newif\ifIncludeLegalRegisters
\IncludeLegalRegistersfalse
\newif\ifIncludeAIRegister
\IncludeAIRegistertrue
\newif\ifIncludeAppendix
\IncludeAppendixtrue

% ---------------------------------------------------------------------------
% Sprache, Schriften und Seitenlayout
% ---------------------------------------------------------------------------
\usepackage[T1]{fontenc}
\usepackage[utf8]{inputenc}
\usepackage[english,ngerman]{babel}
\usepackage{newtxtext,newtxmath}
\usepackage{microtype}
\usepackage{xcolor}
\usepackage{csquotes}
\usepackage{etoolbox}

\usepackage[
  a4paper,
  portrait,
  left=3cm,
  right=4cm,
  top=2.5cm,
  bottom=2.5cm
]{geometry}

\usepackage{setspace}
\onehalfspacing
\setlength{\parindent}{0pt}
\setlength{\parskip}{6pt}

% Deutsch ist voreingestellt. Babel stellt auch englische Trennung und
% Bezeichnungen bereit.
\AtBeginDocument{\selectlanguage{\MainLanguage}}
\newcommand{\DEEN}[2]{\iflanguage{english}{#2}{#1}}

% Standardmaessig werden drei Ebenen nummeriert: 1, 1.1 und 1.1.1. Fuer eine
% umfangreiche Arbeit beide Werte auf 3 setzen, um Ebene vier zu aktivieren.
\setcounter{secnumdepth}{2}
\setcounter{tocdepth}{2}

% Hauptueberschriften sind 14 pt gross, tiefere Ebenen 12 pt und fett.
\setkomafont{disposition}{\normalfont\bfseries\setstretch{1}}
\setkomafont{chapter}{\fontsize{14}{17}\selectfont}
\setkomafont{section}{\fontsize{12}{14.5}\selectfont}
\setkomafont{subsection}{\fontsize{12}{14.5}\selectfont}
\setkomafont{subsubsection}{\fontsize{12}{14.5}\selectfont}
\RedeclareSectionCommand[beforeskip=0pt,afterskip=18pt]{chapter}
\RedeclareSectionCommand[beforeskip=18pt,afterskip=12pt]{section}
\RedeclareSectionCommand[beforeskip=18pt,afterskip=12pt]{subsection}
\RedeclareSectionCommand[beforeskip=18pt,afterskip=12pt]{subsubsection}

% ---------------------------------------------------------------------------
% Kopfzeile, Fusszeile, Fussnoten, Gleitumgebungen und Quellen
% ---------------------------------------------------------------------------
\usepackage[automark,singlespacing=true]{scrlayer-scrpage}
\clearpairofpagestyles
\ihead{\AuthorName}
\ohead{\MatriculationNumber}
\cfoot{}
\AtBeginDocument{%
  \ifdefstring{\MainLanguage}{english}
    {\cfoot{Page \thepage\ of \pageref*{LastNumberedPage}}}
    {\cfoot{Seite \thepage\ von \pageref*{LastNumberedPage}}}}
\setkomafont{pageheadfoot}{\normalfont\fontsize{10}{12}\selectfont}
\setlength{\headheight}{15pt}
\pagestyle{scrheadings}
\renewcommand*{\chapterpagestyle}{scrheadings}

\usepackage{chngcntr}
\counterwithout{footnote}{chapter}
\setlength{\footnotesep}{16pt}
\deffootnote[0.5cm]{0.5cm}{0cm}{%
  \raggedright\makebox[0.5cm][l]{\textsuperscript{\thefootnotemark}}}

\usepackage{graphicx}
\usepackage{booktabs}
\usepackage{tabularx}
\usepackage{array}
\usepackage{caption}
\captionsetup{
  format=plain,
  justification=justified,
  singlelinecheck=false,
  font={footnotesize,it},
  labelfont={bf,it},
  labelsep=colon,
  skip=6pt
}
\counterwithout{figure}{chapter}
\counterwithout{table}{chapter}

% Diese Umgebung setzt Tabellentext in empfohlenen 10 pt und einzeilig.
\newenvironment{hfhtable}[1][htbp]
  {\begin{table}[#1]\centering\footnotesize\setstretch{1}}
  {\end{table}}
\newcommand{\source}[1]{%
  \par\vspace{3pt}{\footnotesize\itshape
    \DEEN{Quelle}{Source}: #1\par}}
\newcolumntype{Y}{>{\raggedright\arraybackslash}X}

\usepackage[round,authoryear]{natbib}
\setlength{\bibhang}{0.5cm}
\setlength{\bibsep}{6pt}
\AtBeginEnvironment{thebibliography}{\singlespacing}

\usepackage{xurl}
\usepackage[hidelinks]{hyperref}
\usepackage{bookmark}
\usepackage{pdfpages}
\hypersetup{
  pdftitle={\ThesisTitle},
  pdfauthor={\AuthorName},
  pdfsubject={\ThesisType}
}

\newcommand{\TemplateHint}[1]{%
  \par\begingroup\small\itshape\color{gray}#1\par\endgroup}

% Das Titelblatt ist Seite 1, zeigt aber keine Seitenzahl. Bei einer
% anmeldepflichtigen Arbeit ist es durch das offizielle Titelblatt des
% Pruefungsamts zu ersetzen. Der alternative Befehl nimmt eine einseitige PDF.
\newcommand{\MakeThesisTitle}{%
  \setcounter{page}{1}
  \thispagestyle{empty}
  \begin{center}
    \vspace*{1.2cm}
    {\large Hamburger Fern-Hochschule\par}
    \vspace{0.3cm}
    {\Department\par}

    \vfill
    {\large\ThesisType\par}
    \vspace{1.2cm}
    {\bfseries\fontsize{20}{24}\selectfont\ThesisTitle\par}
    \vspace{0.5cm}
    {\large\ThesisSubtitle\par}
    \vfill

    \begin{tabular}{@{}ll@{}}
      \DEEN{Verfasser/in}{Author}: & \AuthorName \\
      \DEEN{Matrikelnummer}{Student ID}: & \MatriculationNumber \\
      \DEEN{Studiengang}{Degree programme}: & \DegreeProgramme \\
      \DEEN{Erstprüfer/in}{First examiner}: & \FirstExaminer \\
      \DEEN{Zweitprüfer/in}{Second examiner}: & \SecondExaminer \\
      \DEEN{Abgabedatum}{Submission date}: & \SubmissionDate
    \end{tabular}
  \end{center}
  \clearpage
}

\newcommand{\MakeOfficialTitlePage}[1]{%
  \includepdf[pages=1,pagecommand={\thispagestyle{empty}}]{#1}
  \setcounter{page}{2}
}

% Die HFH-Word-Datei fuehrt das Inhaltsverzeichnis in sich selbst auf. Dieser
% Befehl bildet das ohne Nummerierung der Ueberschrift nach.
\makeatletter
\newcommand{\ThesisTableOfContents}{%
  \chapter*{\contentsname}
  \phantomsection
  \addcontentsline{toc}{chapter}{\contentsname}
  \@starttoc{toc}
}
\makeatother

\begin{document}

\MakeThesisTitle
% Alternative fuer ein offizielles Titelblatt:
% \MakeOfficialTitlePage{offizielles-titelblatt.pdf}

\ifIncludeAbstract
  \addchap{Abstract}
  \TemplateHint{\DEEN
    {Dieser Abschnitt ist nur für Masterarbeiten vorgesehen. Für eine
     Bachelor-, Seminar-, Haus- oder Projektarbeit ist er einschließlich des
     Befehls \texttt{\string\addchap\{Abstract\}} zu löschen. Bei einer
     Masterarbeit umfasst der Abstract etwa eine, höchstens zwei Seiten und
     fasst Forschungsanlass, Fragestellung, theoretische Basis, Methode,
     Ergebnisse und Schlussfolgerung eigenständig verständlich zusammen.}
    {This section is intended for master's theses only. Delete it for a
     bachelor's thesis, seminar paper, term paper, or project paper. A
     master's abstract should be about one and no more than two pages and
     concisely cover the research context, question, theory, method, findings,
     and conclusion.}}

  \DEEN
    {Hier steht der kurze, ergebnisorientierte Abstract der Masterarbeit.}
    {Insert the concise, results-oriented master's abstract here.}
\fi

\ThesisTableOfContents
\listoffigures
\listoftables

\addchap{\DEEN{Abkürzungsverzeichnis}{List of Abbreviations}}
\begin{tabular}{@{}p{2.5cm}%
  p{\dimexpr\textwidth-2.5cm-2\tabcolsep\relax}@{}}
  HFH & Hamburger Fern-Hochschule \\
  KI  & Künstliche Intelligenz
\end{tabular}
\TemplateHint{\DEEN
  {Einträge alphabetisch sortieren. Erklärungsbedürftige Abkürzungen werden
   zusätzlich bei der ersten Verwendung im Text ausgeschrieben.}
  {Sort entries alphabetically. In addition, spell out each non-standard
   abbreviation at first use in the text.}}

% ---------------------------------------------------------------------------
% Haupttext
% ---------------------------------------------------------------------------
\chapter{\DEEN{Einleitung}{Introduction}}

\section{\DEEN{Problemstellung und Zielsetzung}{Problem and Objective}}

Die Hamburger Fern-Hochschule (HFH) bildet den institutionellen Rahmen dieser
Arbeit. Hier werden Problemstellung, Forschungsfrage, Zielsetzung und Aufbau
präzise eingeführt. Wissenschaftliche Aussagen werden mit geeigneten Quellen
belegt \citep[vgl.][S.~21--24]{booth2016}.

Ein neuer Gedanke beginnt in einem neuen Absatz. Der Text ist im Blocksatz
gesetzt und wird automatisch getrennt. Eine Fußnote demonstriert die
vorgesehene Form.\footnote{Diese beispielhafte Fußnote beginnt mit einem
Großbuchstaben und endet mit einem Punkt.}

\section{\DEEN{Aufbau der Arbeit}{Structure of the Thesis}}

Beschreiben Sie knapp, wie die folgenden Kapitel zur Beantwortung der
Forschungsfrage beitragen. Ein Gliederungspunkt sollte nie nur einen einzigen
Unterpunkt besitzen.

\chapter{\DEEN{Theoretischer Rahmen}{Theoretical Framework}}

\section{\DEEN{Begriffe und Modelle}{Concepts and Models}}

Definieren und ordnen Sie die zentralen Begriffe und Modelle ein. Die
Gliederung sollte in der Regel nicht tiefer als drei Ebenen reichen.

\subsection{\DEEN{Zentraler Begriff}{Central Concept}}

Entwickeln Sie den theoretischen Bezugsrahmen aus der Fachliteratur.

\subsection{\DEEN{Weiterführendes Modell}{Additional Model}}

Grenzen Sie konkurrierende Ansätze voneinander ab und begründen Sie Ihre
Auswahl.

\chapter{\DEEN{Methodik}{Methodology}}

\section{\DEEN{Forschungsdesign}{Research Design}}

Dokumentieren Sie das Vorgehen so, dass es fachlich nachvollziehbar ist.
Formeln werden mit kapitelbezogener Nummerierung gesetzt:
\begin{equation}
  \bar{x} = \frac{1}{n}\sum_{i=1}^{n} x_i.
  \label{eq:mean}
\end{equation}
Gleichung~\eqref{eq:mean} wird im laufenden Text erläutert.

\section{\DEEN{Datengrundlage}{Data Basis}}

Tabellen und Abbildungen werden im Text ausdrücklich aufgegriffen. Tabelle
\ref{tab:sample} zeigt eine kompakte Ergebnisdarstellung.

\begin{hfhtable}
  \caption{\DEEN{Beispielhafte Ergebnisübersicht}{Example Results Overview}}
  \label{tab:sample}
  \begin{tabularx}{\textwidth}{@{}Yrr@{}}
    \toprule
    \textbf{\DEEN{Kategorie}{Category}} &
    \textbf{\DEEN{Anzahl}{Count}} &
    \textbf{\DEEN{Anteil}{Share}} \\
    \midrule
    A & 24 & 60\,\% \\
    B & 16 & 40\,\% \\
    \bottomrule
  \end{tabularx}
  \source{\DEEN{Eigene Darstellung}{Author's own work}.}
\end{hfhtable}

\chapter{\DEEN{Ergebnisse und Diskussion}{Results and Discussion}}

\section{\DEEN{Ergebnisse}{Results}}

Abbildung~\ref{fig:sample} dient als Platzhalter. Jede Abbildung muss ohne
unnötige Rückfragen verständlich beschriftet, fortlaufend nummeriert und im
Text referenziert sein.

\begin{figure}[htbp]
  \centering
  \fbox{\rule{0pt}{4cm}\rule{0.72\textwidth}{0pt}}
  \caption{\DEEN{Platzhalter für eine Abbildung}{Figure Placeholder}}
  \label{fig:sample}
  \source{\DEEN{Eigene Darstellung}{Author's own work}.}
\end{figure}

\section{\DEEN{Diskussion}{Discussion}}

Ordnen Sie die Ergebnisse in den Forschungsstand ein, diskutieren Sie ihre
Tragweite und benennen Sie Grenzen der Untersuchung.

\chapter{\DEEN{Fazit und Ausblick}{Conclusion and Outlook}}

Beantworten Sie die Forschungsfrage auf Grundlage der Ergebnisse. Trennen Sie
Schlussfolgerungen, Limitationen und weiterführenden Forschungsbedarf klar.

% ---------------------------------------------------------------------------
% Verzeichnisse nach dem Haupttext
% ---------------------------------------------------------------------------
\renewcommand{\bibname}{\DEEN{Quellenverzeichnis}{References}}
\bibliographystyle{plainnat}
\bibliography{references}

\ifIncludeLegalRegisters
  \addchap{\DEEN{Rechtsprechungsverzeichnis}{Table of Cases}}
  \TemplateHint{\DEEN
    {Nur bei verwendeten Gerichtsentscheidungen behalten. Entscheidungen mit
     Gericht, Entscheidungsart, Datum, Aktenzeichen und Fundstelle aufführen.}
    {Retain only if court decisions were cited. Give the court, decision type,
     date, case number, and publication reference.}}

  \addchap{\DEEN
    {Verzeichnis der Verwaltungsanweisungen und Parlamentaria}
    {Administrative and Parliamentary Sources}}
  \TemplateHint{\DEEN
    {Nur bei Bedarf behalten und die verwendeten amtlichen Dokumente mit
     eindeutiger Fundstelle aufführen.}
    {Retain only when needed and identify every cited official document
     unambiguously.}}
\fi

\ifIncludeAIRegister
  \addchap{\DEEN
    {Verzeichnis der eingesetzten KI-Werkzeuge}
    {Register of AI Tools Used}}
  \TemplateHint{\DEEN
    {Diesen Abschnitt nur behalten, wenn KI-Werkzeuge eingesetzt wurden. Die
     Einträge müssen den tatsächlichen Einsatz nachvollziehbar dokumentieren.}
    {Retain this section only if AI tools were used. Entries must document the
     actual use transparently.}}

  \begingroup
  \footnotesize
  \setstretch{1}
  \begin{tabularx}{\textwidth}{@{}YYYY@{}}
    \toprule
    \textbf{\DEEN{Werkzeug}{Tool}} &
    \textbf{\DEEN{Betroffener Teil}{Affected Part}} &
    \textbf{\DEEN{Art des Einsatzes}{Type of Use}} &
    \textbf{\DEEN{Prüfung und Übernahme}{Review and Adoption}} \\
    \midrule
    \textit{Name und Version} & \textit{Kapitel/Arbeitsschritt} &
    \textit{Funktion} &
    \textit{Kritische Prüfung und konkrete Verwendung} \\
    \bottomrule
  \end{tabularx}
  \endgroup
\fi

\ifIncludeAppendix
  \addchap{\DEEN{Anlagenverzeichnis}{List of Appendices}}
  \textbf{\DEEN{Anlage A}{Appendix A}:}
  \nameref{app:material}\dotfill\pageref{app:material}

  \appendix
  \chapter{\DEEN{Ergänzendes Material}{Supplementary Material}}
  \label{app:material}
  Fügen Sie hier nur Material ein, das für das Verständnis notwendig ist,
  aber den Lesefluss im Haupttext stören würde. Der Anhang ist nicht dazu
  gedacht, die vorgegebene Seitenbegrenzung zu umgehen.
\fi

% Diese Marke bestimmt die Gesamtzahl in der Fusszeile. Die abschliessende
% Erklaerung ist absichtlich ausgenommen und zeigt keine Seitenzahl.
\phantomsection
\label{LastNumberedPage}
\clearpage
\pagestyle{empty}
\chapter*{\DEEN{Eigenständigkeitserklärung}{Declaration of Authorship}}
\thispagestyle{empty}

\TemplateHint{\DEEN
  {Dies ist nur ein Platzhalter. Fügen Sie hier als letzte Seite die aktuelle,
   unveränderte Eigenständigkeitserklärung aus dem WebCampus ein. Sie darf
   weder eine sichtbare Seitenzahl noch einen Eintrag im Inhaltsverzeichnis
   erhalten und muss entsprechend den aktuellen Vorgaben datiert und
   unterschrieben werden.}
  {This is a placeholder only. Insert the current, unmodified declaration
   supplied through WebCampus as the final page. It must have neither a visible
   page number nor a table-of-contents entry and must be dated and signed as
   currently required.}}

% Zum Einfuegen der offiziellen Erklaerung den Platzhalter oben ersetzen durch:
% \includepdf[pages=-,pagecommand={\thispagestyle{empty}}]
%   {eigenstaendigkeitserklaerung.pdf}

\end{document}
